\documentclass[11pt,a4paper]{article}
\usepackage[utf8]{inputenc}
\usepackage[T1]{fontenc}
\usepackage{tgpagella}
\usepackage[spanish]{babel}
\usepackage{geometry}
\usepackage{booktabs}
\usepackage{xcolor}
\usepackage{titlesec}
\usepackage{fancyhdr}
\usepackage[parfill]{parskip}
\usepackage[expansion=false]{microtype}
\usepackage{hyperref}
\usepackage{setspace}
\usepackage{needspace}
\usepackage{etoolbox}
\usepackage{graphicx}

\widowpenalty=10000
\clubpenalty=10000
\displaywidowpenalty=10000
\raggedbottom
\sloppy
\hbadness=10000
\hfuzz=2pt

\geometry{top=3.0cm,bottom=3.0cm,left=3.5cm,right=3.5cm}
\setlength{\parskip}{0.65em}
\setlength{\parindent}{0em}
\setlength{\headheight}{14pt}

\definecolor{azulai}{RGB}{30,80,140}
\definecolor{doradoai}{RGB}{140,90,0}
\definecolor{grisai}{RGB}{70,70,70}

\titleformat{\section}{\Large\bfseries\color{azulai}}{}{0em}{}[{\color{azulai}\titlerule[0.5pt]}]
\titlespacing*{\section}{0pt}{1.8em}{0.8em}

\titleformat{\subsection}{\large\bfseries\color{doradoai}}{}{0em}{}
\titlespacing*{\subsection}{0pt}{1.4em}{0.5em}

\titleformat{\subsubsection}{\normalsize\bfseries\color{grisai}}{}{0em}{}
\titlespacing*{\subsubsection}{0pt}{1.0em}{0.3em}

% Evita títulos huérfanos al final de página
\pretocmd{\section}{\Needspace{12\baselineskip}}{}{}
\pretocmd{\subsection}{\Needspace{9\baselineskip}}{}{}
\pretocmd{\subsubsection}{\Needspace{6\baselineskip}}{}{}

\pagestyle{fancy}
\fancyhf{}
\rhead{\textcolor{grisai}{\small Urko — Arquitectura Interna}}
\lhead{\textcolor{grisai}{\small Luna · Casa 4 · Casa 6}}
\cfoot{\textcolor{grisai}{\small\thepage}}
\renewcommand{\headrulewidth}{0.3pt}

\hypersetup{colorlinks=true,linkcolor=azulai,urlcolor=azulai}
\setstretch{1.45}
\tolerance=2500
\emergencystretch=8em

\begin{document}

% ── Portada ──────────────────────────────────────────────────────────────────
\begin{titlepage}
  \centering
  \vspace*{1.5cm}
  {\Huge\bfseries\color{azulai} Luna · Casa 4 · Casa 6}\\[0.5cm]
  {\large\color{grisai} Arquitectura Interna}\\[0.3cm]
  {\small\itshape\color{grisai} Un método para sostener cuerpo, energía y vida con coherencia}\\[2cm]
  {\huge\color{doradoai} Urko}\\[1.5cm]
  {\Large 11/12/2015 \quad 22:47}\\[0.3cm]
  {\Large Pamplona, España}\\[0.3cm]
  {\normalsize Lat: 42.8157° \quad Lon: -1.6522° \quad UTC+01:00}\\[0.3cm]
  {\normalsize Ascendente: Leo 25°29' \quad
    MC: Tauro 17°51'}\\[2cm]

  \begin{tabular}{lll}
    \textbf{Luna:} & Sagitario & Casa 5 \\
    \textbf{Casa 4:} & Escorpio en cúspide & Regente: Plutón \\
    \textbf{Casa 6:} & Capricornio en cúspide & Regente: Saturno \\
  \end{tabular}\\[2cm]

  \vfill
  {\small Generado el 09/05/2026}
\end{titlepage}

\tableofcontents
\clearpage

% ── Datos de referencia ───────────────────────────────────────────────────────
\section{Datos de referencia}

\begin{center}
\begin{tabular}{llll}
  \toprule
  \textbf{Eje} & \textbf{Signo} & \textbf{Posición} & \textbf{Datos} \\
  \midrule
  Luna & Sagitario & Casa 5 & 25°07' \\
  Casa 4 (cúspide) & Escorpio & --- & Regente: Plutón \\
  Casa 6 (cúspide) & Capricornio & --- & Regente: Saturno \\
  Ascendente & Leo & --- & 25°29' \\
  Medio Cielo & Tauro & --- & 17°51' \\
  \bottomrule
\end{tabular}
\end{center}

\vspace{0.4cm}
\textbf{Planetas en Casa 4:} Sol, Saturno\\
\textbf{Planetas en Casa 6:} ninguno

\vspace{0.5cm}
\textbf{Aspectos de la Luna:}

\begin{center}
\begin{tabular}{lllll}
  \toprule
  & \textbf{Asp.} & \textbf{Planeta} & \textbf{Tipo} & \textbf{Orbe} \\
  \midrule
  Luna & □ & Nodo Norte & Cuadratura & 2.0° \\
  Luna & □ & Nodo Sur & Cuadratura & 2.0° \\
  Luna & □ & Júpiter & Cuadratura & 3.0° \\
  Luna & □ & Quirón & Cuadratura & 5.3° \\
  Luna & = & Sol & Conjunción & 5.6° \\
  Luna & = & Mercurio & Conjunción & 7.6° \\
  \bottomrule
\end{tabular}
\end{center}

\vspace{0.8cm}
\begin{center}
\includegraphics[width=0.80\textwidth]{C:/Users/Izas/Claude/Astrología/Pruebas/Urko_Luna_C4_C6_rueda.png}
\end{center}

\clearpage

% ── Interpretación ────────────────────────────────────────────────────────────
\section{Interpretación — Arquitectura Interna}

\Needspace{8\baselineskip}
\begin{center}
{\small\itshape
No se trata de explicarte. Se trata de mostrar cómo funciona tu regulación interna,\\
dónde puede haber desregulación y qué necesitas para sostenerte.
}
\end{center}

\vspace{0.5cm}

% ── 1. Luna ───────────────────────────────────────────────────────────────────
\subsection{1. Tu Luna — Función emocional}

\subsubsection*{Luna en Sagitario — Casa 5}

Necesitas sentir que lo que estás viviendo tiene dirección.

Cuando una emoción encuentra sentido o movimiento, puedes atravesarla mucho mejor. El problema aparece cuando sientes algo que no sabes hacia dónde va o que parece no conducir a ningún lugar.

Entonces suele aparecer inquietud, necesidad de escapar, cambiar de aire o salir de la situación de alguna manera. Tu cuerpo necesita espacio, movimiento y sensación de horizonte.

Cuando pasas demasiado tiempo sintiendo que no existe libertad o amplitud suficiente, la activación empieza a acumularse.

La regulación mejora cuando existe dirección, expansión y posibilidad de avanzar. En cambio, la sensación de límite, encierro o falta de sentido tiende a aumentar rápidamente la inquietud interna.

Necesitas expresar lo que sientes para mantener vitalidad emocional.

La creatividad, el juego, el disfrute o la posibilidad de ser espontáneo funcionan como formas reales de regulación interna.

Cuando pasas demasiado tiempo reprimiendo esa parte de ti, la energía emocional empieza a acumularse y puede aparecer irritabilidad, vacío o sensación de desconexión contigo.

Tu cuerpo suele apagarse cuando la vida pierde espacio para el disfrute auténtico o para la expresión personal.

La regulación mejora cuando puedes crear, disfrutar y expresarte libremente. En cambio, la exigencia constante, la represión emocional o sentir que no tienes espacio para ser quien eres tienden a desgastar profundamente la energía.

Tu mundo emocional puede entrar en tensión con la dirección de crecimiento de tu vida.

A veces lo que resulta familiar emocionalmente no es lo que más ayuda a avanzar. El crecimiento exige reconocer patrones antiguos sin seguir obedeciéndolos automáticamente.

La regulación aparece cuando puedes distinguir entre seguridad conocida y evolución real.

Hay patrones emocionales antiguos que pueden generar tensión interna.

A veces reaccionas desde una memoria emocional que ya no corresponde del todo al presente. Algo se activa ahora, pero la respuesta puede venir de un lugar anterior.

El trabajo está en reconocer cuándo algo pertenece al momento actual y cuándo pertenece a una forma antigua de protección.

Muchas veces puedes sentir más necesidad emocional de la que el entorno puede cubrir fácilmente.

Cuando algo te hace sentir bien, quieres más. Y cuando falta, el vacío puede sentirse enorme. Internamente cuesta encontrar un punto de suficiente.

La regulación necesita medida, no negación de la necesidad.

Hay una herida emocional que puede activarse en momentos de necesidad, cuidado o pertenencia.

A veces lo que más necesitas también es lo que más cuesta recibir. La sensibilidad se activa justo en el lugar donde haría falta descanso o apoyo.

El trabajo no es corregir esa sensibilidad, sino aprender a acompañarla sin rechazarla.

Lo que sientes y la dirección vital suelen moverse muy cerca.

Cuando emocionalmente hay estabilidad, aparece más energía, claridad y sensación de propósito. Y cuando algo se rompe por dentro, la fuerza vital también lo nota enseguida.

No separas fácilmente quién eres de cómo te sientes en ese momento. Eso puede darte mucha coherencia interna, pero también hacer que los períodos difíciles afecten con más profundidad.

Pensar y sentir ocurren casi al mismo tiempo.

Necesitas poner palabras a lo que te pasa para terminar de integrarlo. Cuando algo no puede expresarse, la emoción sigue girando dentro de la mente y cuesta más que encuentre reposo.

Hablar, escribir o compartir lo que sientes puede ayudarte mucho a recuperar claridad interna.

% ── 2. Casa 4 ─────────────────────────────────────────────────────────────────
\subsection{2. Tu Casa 4 — Raíz interna}

\subsubsection*{Escorpio en la cúspide — Regente: Plutón}

Percibes muy rápido lo que ocurre debajo de la superficie.

Las tensiones no dichas, las emociones contenidas o las dinámicas ocultas te afectan incluso cuando nadie las nombra.

Eso hace que muchas veces permanezcas en vigilancia aunque externamente todo parezca tranquilo.

Necesitas sentir confianza profunda para relajarte de verdad. Cuando percibes manipulación, secretos o falta de honestidad emocional, el cuerpo vuelve rápidamente al estado de alerta.

La regulación mejora cuando existe intimidad real, honestidad emocional y confianza profunda. En cambio, la manipulación, las tensiones ocultas o la sensación de amenaza emocional tienden a activar mucha intensidad interna.

El regente de la Casa 4 es Plutón, en Capricornio y en Casa 5. La base interna tiende a organizarse desde la creatividad y la expresión, con la cualidad concreta y sostenida de Capricornio.

Necesitas sentir que puedes ser plenamente quien eres dentro de tu espacio más íntimo.

Cuando existe demasiada distancia entre lo que muestras hacia fuera y lo que realmente sucede dentro de ti, algo empieza a desgastarse lentamente.

No suele regularte sostener una imagen que no tiene raíz en tu vida real. Tu mundo interno necesita coherencia.

El hogar no funciona solo como lugar físico. Es el espacio donde deberías poder bajar la exigencia y existir sin tener que sostener una actuación constante.

Cuando no existe ese lugar de autenticidad y descanso interno, la energía empieza a agotarse silenciosamente.

Aprendiste pronto a sostenerte emocionalmente sin depender demasiado de otras personas.

Puede que en el entorno de origen no hubiera demasiado espacio para mostrar vulnerabilidad libremente.

Eso generó fortaleza. Y también cierta dificultad para relajarte completamente o recibir apoyo sin medirlo antes.

Muchas veces aparece la sensación de tener que merecer el descanso, el cuidado o la ayuda.

Tu seguridad interna se construye lentamente, a través de tiempo, estructura y experiencia real.

% ── 3. Casa 6 ─────────────────────────────────────────────────────────────────
\subsection{3. Tu Casa 6 — Regulación cotidiana}

\subsubsection*{Capricornio en la cúspide — Regente: Saturno}

Puedes sostener muchísimo más de lo que la mayoría imagina.

La dificultad es que también puedes ignorar las señales de cansancio durante demasiado tiempo.

Existe tendencia a seguir funcionando incluso cuando el cuerpo ya está pidiendo parar.

Muchas veces las responsabilidades aparecen antes que la recuperación.

No porque no sientas el agotamiento, sino porque el umbral interno para detenerte suele ser muy alto.

La desregulación suele aparecer de golpe, cuando la acumulación ya supera lo que podías sostener.

La regulación mejora cuando existen límites claros, descanso suficiente y estructura estable. En cambio, la sobrecarga prolongada, el exceso de responsabilidad o no permitirte parar tienden a desgastar profundamente el cuerpo.

El regente de la Casa 6 es Saturno, en Sagitario y en Casa 4. La regulación cotidiana tiende a organizarse desde las raíces y el hogar interno, con la cualidad activa y directa de Sagitario.

% ── 4. Integración ────────────────────────────────────────────────────────────
\subsection{4. Integración — Cómo funcionan las tres capas juntas}

Tu Luna en Sagitario muestra cómo reaccionas emocionalmente cuando algo te afecta.

Es el primer movimiento: lo que aparece antes de que puedas ordenarlo, explicarlo o decidir qué hacer con ello.

Después, esa emoción busca un lugar donde apoyarse. Ahí entra tu Casa 4, con Escorpio en la cúspide: la parte de ti que necesita sentir base, refugio y seguridad interna.

Si esa base está disponible, lo que sientes puede moverse con más facilidad.
Si no lo está, la emoción se queda más tiempo en estado de alerta.

Y lo que no encuentra apoyo interno termina llegando al cuerpo. Ahí entra tu Casa 6, con Capricornio en la cúspide: tus hábitos, tu energía diaria, tu descanso y la forma en que el cuerpo sostiene lo que vives.

Por eso, cuando algo se desordena emocionalmente, no se queda solo en lo emocional. Puede acabar apareciendo como cansancio, pérdida de ritmo, dificultad para descansar o sensación de no poder sostener el día con la misma energía.

En tu caso, las tres capas tienen elementos distintos: Fuego en la Luna, Agua en la Casa 4 y Tierra en la Casa 6.

Esto significa que tu emoción, tu base interna y tu cuerpo no siempre piden lo mismo.

Lo que te ayuda emocionalmente puede no ser lo que tu cuerpo necesita ese día. Lo que te da seguridad interna puede no coincidir con el ritmo que sostiene tu rutina.

Por eso es importante no buscar una única respuesta para todo. A veces una parte de ti necesita expresión, otra necesita refugio y otra necesita descanso, orden o movimiento.

Además, el Sol y Saturno en la Casa 4 influye directamente en tu sensación de base interna.

Cuando buscas seguridad, refugio o estabilidad emocional profunda, esa energía ya está presente ahí.

Por eso tu forma de sentir hogar, intimidad y raíz no depende solo del signo de la Casa 4. También está marcada por lo que esos planetas activan dentro de ti.

Los aspectos de tensión con el Nodo Norte, el Nodo Sur y Júpiter hacen que el mundo emocional tenga menos margen en ciertos momentos.

Cuando algo te activa, puede haber menos tiempo entre sentir y reaccionar. La emoción puede llegar más intensa, más mezclada o más difícil de ordenar.

Esto no significa que estés condenado a desregularte. Significa que necesitas reconocer antes las señales iniciales, porque cuando la activación ya ha subido mucho, cuesta más volver al centro.

Al mismo tiempo, los aspectos de apoyo con el Sol y Mercurio ofrecen recursos reales.

Cuando consigues encontrar un primer punto de estabilidad, esos aspectos ayudan a sostenerlo.

No eliminan la tensión, pero pueden facilitar que vuelvas a organizarte sin quedarte completamente dentro de lo que se activó.

Lo más difícil no suele ser que una sola parte se altere.

Lo más difícil aparece cuando las tres capas se mueven al mismo tiempo:
la Luna en Sagitario, con su forma concreta de sentir;
la Casa 4 en Escorpio, con su necesidad de base y seguridad interna;
y la Casa 6 en Capricornio, con el cuerpo intentando sostenerlo todo en la vida diaria.

Cuando esto ocurre, puedes notar que no sabes por dónde empezar: emocionalmente hay activación, internamente falta suelo y el cuerpo empieza a mostrar cansancio, tensión o pérdida de ritmo.

El trabajo no consiste en resolverlo todo a la vez.

Consiste en encontrar la primera capa disponible: una emoción que pueda expresarse, un entorno que pueda darte más seguridad, o una acción corporal sencilla que ayude al sistema a volver poco a poco.

% ── 5. Sostén y desregulación ─────────────────────────────────────────────────
\subsection{5. Sostén y desregulación}

Tu Luna en Sagitario se desregula cuando lo que sientes no encuentra salida.

Puede ser porque no puedes expresarlo, porque no hay espacio para moverte, porque el entorno no lo recibe o porque intentas sostenerlo demasiado tiempo dentro.

No siempre hace falta que ocurra algo muy intenso.

A veces basta con acumular durante demasiado tiempo una emoción que no ha tenido lugar, palabra, movimiento o descanso.

Los aspectos de tensión con el Nodo Norte, el Nodo Sur y Júpiter hacen que ese margen sea más estrecho.

Cuando algo te activa, puede haber menos tiempo entre sentir y reaccionar.

La emoción puede llegar más fuerte, más mezclada o más difícil de ordenar. Por eso es importante reconocer las primeras señales antes de llegar al punto de desborde.

Tu Casa 4 con Escorpio pierde estabilidad cuando dejas de sentir base interna.

No siempre tiene que haber un conflicto claro.

A veces basta con que cambie el clima emocional del entorno, con que una relación cercana se tense, con que no puedas descansar de verdad o con que el espacio íntimo deje de sentirse seguro.

Cuando eso ocurre, no solo cambia el ánimo. Cambia la sensación de tener suelo por dentro.

Tu Casa 6 con Capricornio muestra la desregulación en el cuerpo y en la vida cotidiana.

Puede aparecer como cansancio, pérdida de ritmo, tensión física, sueño menos reparador o dificultad para sostener hábitos que normalmente sí puedes sostener.

Muchas veces el cuerpo avisa antes de que puedas explicar qué está pasando.

Cuando el cuerpo empieza a perder continuidad, conviene mirar también qué emoción no ha encontrado salida y qué parte de tu base interna ha dejado de sentirse segura.

Los momentos más difíciles aparecen cuando las tres capas se mueven a la vez.

La Luna en Sagitario se activa.
La Casa 4 en Escorpio deja de sentirse como base.
Y la Casa 6 en Capricornio empieza a mostrarlo en el cuerpo, el descanso o la rutina.

Ahí puedes sentir que no sabes por dónde empezar.

No se trata de resolverlo todo a la vez.
Se trata de encontrar la primera capa disponible: expresar algo, descansar, ordenar una pequeña rutina, moverte, pedir espacio o recuperar una mínima sensación de seguridad.

Los aspectos de apoyo con el Sol y Mercurio funcionan como recursos cuando consigues recuperar un primer punto de estabilidad.

No eliminan la desregulación, pero pueden ayudarte a reorganizarte antes.

Cuando aparece una mínima base, esos aspectos facilitan que no te quedes dentro de lo que se activó.

% ── 6. Orientación final ──────────────────────────────────────────────────────
\subsection{6. Orientación final}

\subsubsection*{Qué observar}

Empieza observando qué situaciones activan tu Luna en Sagitario.

No para corregirte.
No para controlar lo que sientes.

Solo para reconocer antes cómo empieza la activación.

¿Qué ocurre justo antes de sentirte sobrepasado?
¿Qué cambia en el entorno, en los vínculos o en el ritmo diario?

Muchas veces la desregulación no empieza en el momento del desborde. Empieza bastante antes: cuando tu Casa 4 en Escorpio deja de sentirse segura o cuando tu Casa 6 en Capricornio pierde continuidad y el cuerpo empieza a acumular tensión.

\subsubsection*{Qué puede desregularte}

En tu caso, la desregulación suele seguir una secuencia.

Primero, tu Luna en Sagitario busca algo que necesita para estabilizarse y no lo encuentra.

Después, tu Casa 4 deja de funcionar como base interna. Puede ser por tensión emocional, por cambios en el entorno, por exceso de carga o simplemente porque llevas demasiado tiempo sosteniendo más de lo que puedes absorber.

Y finalmente, la Casa 6 en Capricornio empieza a mostrarlo en el cuerpo:
menos energía,
menos descanso,
menos capacidad de sostener hábitos o continuidad.

Cuando las tres capas coinciden, puedes sentir que todo cuesta más de lo habitual aunque desde fuera no parezca estar ocurriendo nada tan grave.

\subsubsection*{Qué te estructura y sostiene}

Lo que te ayuda a regularte no tiene que ser enorme.

De hecho, muchas veces lo más importante son cosas pequeñas que se sostienen en el tiempo.

Una mínima rutina corporal.
Un espacio de descanso real.
Un entorno que se sienta más seguro.
Una forma de expresar lo que ocurre antes de acumularlo demasiado.

Cuando tu Casa 4 siente suelo, tu Luna necesita menos esfuerzo para estabilizarse.

Y cuando tu Casa 6 mantiene cierta continuidad —aunque sea pequeña— el cuerpo puede absorber mejor lo que emocionalmente estás atravesando.

No se trata de no sentir.
Se trata de que lo que sientes no tenga que terminar descargándose siempre en agotamiento, tensión o pérdida de equilibrio.

\subsubsection*{Práctica}

Tu regulación necesita repetición y continuidad.

Elige una práctica corporal o cotidiana sencilla y mantenla estable en el tiempo.

Puede ser desayunar sin prisa, salir unos minutos al exterior, tener una rutina de descanso o respirar conscientemente siempre en el mismo momento del día.

Lo importante no es la intensidad.
Es la constancia.

La repetición le transmite seguridad al sistema nervioso.

\Needspace{8\baselineskip}
\vspace{0.6cm}

\begin{center}
{\small\itshape\color{grisai}
La astrología se usa aquí como lenguaje simbólico de observación, no como definición de la persona.\\
Lo que tienes delante es un mapa de posibilidades, no un destino fijo.
}
\end{center}

\end{document}
